% GENERATED FILE. Edit canonical lessons, then run npm run generate:books.
% canonical-source-hash: fnv1a64:735de8c5a1457844
% canonical-lessons: BN-C514-photano, BN-C514-tatano, BN-C514-seka, BN-C514-bhapano, BN-C514-satlano

\chapter{At the Stove}
\label{ch:bn-at-the-stove}
\hlchaptermodality{514}

\begin{chapteropening}
\textbf{By the end of this chapter:} \emph{I can say how food is boiled, heated, toasted, steamed and sautéed.}

Complete the last lesson of chapter 514: I can say how food is boiled, heated, toasted, steamed and sautéed.
\end{chapteropening}

\section[\texorpdfstring{\bn{ফোটানো} (phoṭāno)}{ফোটানো (phoṭāno)}]{\bn{ফোটানো} (phoṭāno) — to bring to the boil}

\label{lesson:BN-C514-photano}



\begin{quote}
Before the new one: say the Bengali for to knead, then the Bengali for to chop finely.
\end{quote}

\subsection*{You'll want to know: \bn{ফোটানো}}
\textbf{\bn{ফোটানো}} — \emph{phoṭāno} — \textquotedblleft{}to bring to the boil\textquotedblright{}.

\begin{quote}
\bn{আমি} \bn{চায়ের} \bn{জন্য} \bn{জল} \bn{ফোটালাম।}
\end{quote}

\begin{quote}
\emph{Āmi chāyer jônyô jôl phoṭālām.}
\end{quote}

\begin{quote}
I boiled water for tea.
\end{quote}

Past (he/she …): \bn{তিনি} \bn{দুধ} \bn{ফোটালেন}, \emph{tini dudh phoṭālen}.

\begin{cousinweb}
\textbf{\bn{ফোটানো}} (\emph{phoṭāno}) is the causative of \bn{ফোটা}, \textquotedblleft{}to burst open, to boil\textquotedblright{}, from a reconstructed \emph{sphuṭyati} (CDIAL 13845): making the water burst into bubbles.
\end{cousinweb}

\begin{grammarlens}[title={What you've built}]
One more word for describing what you lived and felt.
\end{grammarlens}

\subsection*{Your turn}
\begin{itemize}
\raggedright
  \item \emph{Say it:} \emph{phoṭāno}
  \item \emph{Say it:} \emph{phoṭāno}, once more
  \item \emph{Say it:} \emph{Āmi chāyer jônyô jôl phoṭālām.} — I boiled water for tea.
  \item \emph{Recall:} say the Bengali for to knead, then the Bengali for to chop finely, then say \emph{phoṭāno} again
\end{itemize}

\begin{tcolorbox}[breakable,colback=teal!4,colframe=teal!35!black,title={Before you move on}]
What does \bn{ফোটানো} mean? (\textquotedblleft{}To bring to the boil\textquotedblright{}.) Say it once more.
\end{tcolorbox}
\section[\texorpdfstring{\bn{তাতানো} (tātāno)}{তাতানো (tātāno)}]{\bn{তাতানো} (tātāno) — to heat up (oil, a pan)}

\label{lesson:BN-C514-tatano}



\begin{quote}
Before the new one: say the Bengali for to chop finely, then the Bengali for to bring to the boil.
\end{quote}

\subsection*{You'll want to know: \bn{তাতানো}}
\textbf{\bn{তাতানো}} — \emph{tātāno} — \textquotedblleft{}to heat up (oil, a pan)\textquotedblright{}.

\begin{quote}
\bn{মাছ} \bn{ভাজার} \bn{আগে} \bn{আমি} \bn{তেল} \bn{তাতালাম।}
\end{quote}

\begin{quote}
\emph{Māchh bhājār āge āmi tel tātālām.}
\end{quote}

\begin{quote}
I heated the oil before frying the fish.
\end{quote}

Past (he/she …): \bn{তিনি} \bn{কড়াই} \bn{তাতালেন}, \emph{tini kôṛāi tātālen}.

\begin{cousinweb}
\textbf{\bn{তাতানো}} (\emph{tātāno}) is a tadbhava: CDIAL lists Bengali \emph{tātāna}, \textquotedblleft{}to heat\textquotedblright{}, under Sanskrit \emph{tapta}, \textquotedblleft{}heated, hot\textquotedblright{} (CDIAL 5679).
\end{cousinweb}

\begin{grammarlens}[title={What you've built}]
One more word for describing what you lived and felt.
\end{grammarlens}

\subsection*{Your turn}
\begin{itemize}
\raggedright
  \item \emph{Say it:} \emph{tātāno}
  \item \emph{Say it:} \emph{tātāno}, once more
  \item \emph{Say it:} \emph{Māchh bhājār āge āmi tel tātālām.} — I heated the oil before frying the fish.
  \item \emph{Recall:} say the Bengali for to chop finely, then the Bengali for to bring to the boil, then say \emph{tātāno} again
\end{itemize}

\begin{tcolorbox}[breakable,colback=teal!4,colframe=teal!35!black,title={Before you move on}]
What does \bn{তাতানো} mean? (\textquotedblleft{}To heat up (oil, a pan)\textquotedblright{}.) Say it once more.
\end{tcolorbox}
\section[\texorpdfstring{\bn{সেঁকা} (sẽkā)}{সেঁকা (sẽkā)}]{\bn{সেঁকা} (sẽkā) — to toast; to roast lightly}

\label{lesson:BN-C514-seka}



\begin{quote}
Before the new one: say the Bengali for to bring to the boil, then the Bengali for to heat up.
\end{quote}

\subsection*{You'll want to know: \bn{সেঁকা}}
\textbf{\bn{সেঁকা}} — \emph{sẽkā} — \textquotedblleft{}to toast; to roast lightly\textquotedblright{}.

\begin{quote}
\bn{আমি} \bn{সকালে} \bn{রুটি} \bn{সেঁকলাম।}
\end{quote}

\begin{quote}
\emph{Āmi sôkāle ruṭi sẽklām.}
\end{quote}

\begin{quote}
I toasted bread in the morning.
\end{quote}

Past (he/she …): \bn{তিনি} \bn{রুটি} \bn{সেঁকলেন}, \emph{tini ruṭi sẽklen}.

\begin{cousinweb}
\textbf{\bn{সেঁকা}} (\emph{sẽkā}) comes from a reconstructed root \emph{sēkk}; CDIAL lists Bengali \emph{sẽkā}, \textquotedblleft{}to roast, to sear\textquotedblright{} (CDIAL 13581).
\end{cousinweb}

\begin{grammarlens}[title={What you've built}]
One more word for describing what you lived and felt.
\end{grammarlens}

\subsection*{Your turn}
\begin{itemize}
\raggedright
  \item \emph{Say it:} \emph{sẽkā}
  \item \emph{Say it:} \emph{sẽkā}, once more
  \item \emph{Say it:} \emph{Āmi sôkāle ruṭi sẽklām.} — I toasted bread in the morning.
  \item \emph{Recall:} say the Bengali for to bring to the boil, then the Bengali for to heat up, then say \emph{sẽkā} again
\end{itemize}

\begin{tcolorbox}[breakable,colback=teal!4,colframe=teal!35!black,title={Before you move on}]
What does \bn{সেঁকা} mean? (\textquotedblleft{}To toast; to roast lightly\textquotedblright{}.) Say it once more.
\end{tcolorbox}
\section[\texorpdfstring{\bn{ভাপানো} (bhāpāno)}{ভাপানো (bhāpāno)}]{\bn{ভাপানো} (bhāpāno) — to steam (food)}

\label{lesson:BN-C514-bhapano}



\begin{quote}
Before the new one: say the Bengali for to heat up, then the Bengali for to toast.
\end{quote}

\subsection*{You'll want to know: \bn{ভাপানো}}
\textbf{\bn{ভাপানো}} — \emph{bhāpāno} — \textquotedblleft{}to steam (food)\textquotedblright{}.

\begin{quote}
\bn{দিদিমা} \bn{ইলিশ} \bn{মাছ} \bn{ভাপালেন।}
\end{quote}

\begin{quote}
\emph{Didimā ilish māchh bhāpālen.}
\end{quote}

\begin{quote}
My grandmother steamed some hilsa.
\end{quote}

Past (he/she …): \bn{তিনি} \bn{মাছ} \bn{ভাপালেন}, \emph{tini māchh bhāpālen}.

\begin{cousinweb}
\textbf{\bn{ভাপানো}} (\emph{bhāpāno}) is a tadbhava: CDIAL lists Bengali \emph{bhāpā}, \textquotedblleft{}to emit steam\textquotedblright{}, under Sanskrit \emph{bāṣpāyate} (CDIAL 9225).
\end{cousinweb}

\begin{grammarlens}[title={What you've built}]
One more word for describing what you lived and felt.
\end{grammarlens}

\subsection*{Your turn}
\begin{itemize}
\raggedright
  \item \emph{Say it:} \emph{bhāpāno}
  \item \emph{Say it:} \emph{bhāpāno}, once more
  \item \emph{Say it:} \emph{Didimā ilish māchh bhāpālen.} — My grandmother steamed some hilsa.
  \item \emph{Recall:} say the Bengali for to heat up, then the Bengali for to toast, then say \emph{bhāpāno} again
\end{itemize}

\begin{tcolorbox}[breakable,colback=teal!4,colframe=teal!35!black,title={Before you move on}]
What does \bn{ভাপানো} mean? (\textquotedblleft{}To steam (food)\textquotedblright{}.) Say it once more.
\end{tcolorbox}
\section[\texorpdfstring{\bn{সাঁতলানো} (sā̃tlāno)}{সাঁতলানো (sā̃tlāno)}]{\bn{সাঁতলানো} (sā̃tlāno) — to sauté (in a little oil)}

\label{lesson:BN-C514-satlano}



\begin{quote}
Before the new one: say the Bengali for to toast, then the Bengali for to steam.
\end{quote}

\subsection*{You'll want to know: \bn{সাঁতলানো}}
\textbf{\bn{সাঁতলানো}} — \emph{sā̃tlāno} — \textquotedblleft{}to sauté (in a little oil)\textquotedblright{}.

\begin{quote}
\bn{আমি} \bn{পেঁয়াজ} \bn{আর} \bn{রসুন} \bn{সাঁতলালাম।}
\end{quote}

\begin{quote}
\emph{Āmi pẽyāj ār rôsun sā̃tlālām.}
\end{quote}

\begin{quote}
I sautéed the onion and garlic.
\end{quote}

Past (he/she …): \bn{তিনি} \bn{সবজি} \bn{সাঁতলালেন}, \emph{tini sôbji sā̃tlālen}.

\begin{cousinweb}
\textbf{\bn{সাঁতলানো}} (\emph{sā̃tlāno}) comes from a reconstructed \emph{saṁtalati}, \textquotedblleft{}fries\textquotedblright{}; CDIAL lists Bengali \emph{sā̃talā}, \textquotedblleft{}to singe in oil or ghee\textquotedblright{} (CDIAL 12883).
\end{cousinweb}

\begin{grammarlens}[title={What you've built}]
One more word for describing what you lived and felt.
\end{grammarlens}

\subsection*{Your turn}
\begin{itemize}
\raggedright
  \item \emph{Say it:} \emph{sā̃tlāno}
  \item \emph{Say it:} \emph{sā̃tlāno}, once more
  \item \emph{Say it:} \emph{Āmi pẽyāj ār rôsun sā̃tlālām.} — I sautéed the onion and garlic.
  \item \emph{Recall:} say the Bengali for to toast, then the Bengali for to steam, then say \emph{sā̃tlāno} again
\end{itemize}

\begin{tcolorbox}[breakable,colback=teal!4,colframe=teal!35!black,title={Before you move on}]
What does \bn{সাঁতলানো} mean? (\textquotedblleft{}To sauté (in a little oil)\textquotedblright{}.) Say it once more.
\end{tcolorbox}
